% !PRINTED 298
\weberpage{298}
%<<CONT
\[
\begin{pmatrix} 1, & 1 \\ 2, & -2 \end{pmatrix}\begin{pmatrix} 0, & 1 \\ -1, & 0 \end{pmatrix}\begin{pmatrix} -2, & -1 \\ -2, & 1 \end{pmatrix}\begin{pmatrix} 0, & 1 \\ -1, & 0 \end{pmatrix} = \begin{pmatrix} -2, & 0 \\ 0, & 2 \end{pmatrix},
\]
\[
\begin{pmatrix} -2, & 0 \\ 0, & 2 \end{pmatrix}\begin{pmatrix} 1, & 0 \\ -1, & 1 \end{pmatrix}\begin{pmatrix} 2, & 0 \\ 0, & -2 \end{pmatrix}\begin{pmatrix} 1, & 0 \\ 1, & 1 \end{pmatrix} = \begin{pmatrix} 1, & 0 \\ 2, & 1 \end{pmatrix},
\]
i.\,e. $A^2$, and hence all powers of $A$, from which it is to be concluded, as above, that $\mathfrak{N}$ is identical with $\mathfrak{L}_0$. We have therefore the theorem:

\gesperrt{If $p > 3$, then the group $\mathfrak{L}_0$ has no normal divisor.}

In the case $p = 3$, $\mathfrak{N}$ likewise contains
\[
\begin{pmatrix} 0, & 1 \\ -1, & 0 \end{pmatrix},
\]
and therefore also
\[
\begin{pmatrix} 1, & 0 \\ \pm1, & 1 \end{pmatrix}\begin{pmatrix} 0, & 1 \\ -1, & 0 \end{pmatrix}\begin{pmatrix} 1, & 0 \\ \mp1, & 1 \end{pmatrix} = \begin{pmatrix} \mp1, & 1 \\ 1, & \pm1 \end{pmatrix},
\]
and in fact the four substitutions
\weq{15}{\begin{pmatrix} 1, & 0 \\ 0, & 1 \end{pmatrix},\ \begin{pmatrix} 0, & 1 \\ -1, & 0 \end{pmatrix},\ \begin{pmatrix} -1, & 1 \\ 1, & 1 \end{pmatrix},\ \begin{pmatrix} 1, & 1 \\ 1, & -1 \end{pmatrix}}
do form a normal divisor of $\mathfrak{L}_0$ of index 3.

In the case $p = 3$ the group $\mathfrak{L}_0$ is isomorphic with the alternating group of the permutations of four elements. This is the group of an \gesperrt{arbitrary} equation of the fourth degree, when the square root of the discriminant is adjoined to the domain of rationality. The divisor (15) of this group, of index 3, then furnishes the cubic resolvent of the biquadratic equation that is known in algebra.

A function \gesperrt{belonging} to the group (15), formed from the roots $v_\infty, v_0, v_1, v_2$ of the transformation equation for the third degree of transformation, is for example
\[
v_\infty v_0 + v_1 v_2
\]
or
\[
(v_\infty - v_0)(v_1 - v_2).
\]
The difference between these two functions is this, that the first remains unchanged under the substitution of determinant $-1$
\[
\begin{pmatrix} -1, & 0 \\ 0, & 1 \end{pmatrix}
\]
while the second thereby changes its sign. The first will therefore lead to an equation of the third degree with \gesperrt{rational} coefficients, while in the cubic equation for the second $\sqrt{-3}$ still occurs (§ 78).
